%!TEX TS-program = pdflatex
% Python Intro -- Digitale Technologien
% Die Schalter bleiben zur Kompatibilitaet mit der Vorlage erhalten.
\newcommand{\printpraesenzlsg}{false}
\newcommand{\printloesungen}{false}
\newcommand{\printbewertungen}{false}
\newcommand{\blattnummer}{1}
\documentclass[%
  paper=a4,
  fontsize=10pt,
  ngerman
  ]{scrartcl}

% Basics für Codierung und Sprache
% ===========================================================
  \usepackage{scrtime}
  \usepackage{etex}
  \usepackage{shellesc}
  \usepackage[final]{graphicx}
  \usepackage[utf8]{inputenc}
  \usepackage{babel}
  \usepackage[german=quotes]{csquotes}
  \usepackage{datetime2}
  \usepackage{ifthen}
  \usepackage{dingbat}
% ===========================================================

%% Fonts und Typographie
%% ===========================================================
  \usepackage{charter}
  \usepackage{nimbusmononarrow}
  \usepackage[babel=true,final,tracking=smallcaps]{microtype}
  \DisableLigatures{encoding = T1, family = tt* }
  \usepackage{ellipsis}
% ===========================================================

% Farben
% ===========================================================
  \usepackage[x11names,table]{xcolor}
  \definecolor{fbblau}{HTML}{3078AB}
  \definecolor{blackberry}{rgb}{0.53, 0.0, 0.25}
% ===========================================================

% Mathe-Pakete und -Einstellungen
% ===========================================================
  \usepackage{mathtools}
  \usepackage{amssymb}
  \usepackage[bigdelims]{newtxmath}
  \allowdisplaybreaks
  \usepackage{bm}
  \usepackage{wasysym}
  \usepackage{bussproofs}
  % skalierbarer Beweisbaum: \begin{scprooftree}{0.8} ... \end{scprooftree}
  \newenvironment{scprooftree}[1]%
    {\gdef\scalefactor{#1}\center\proofSkipAmount \leavevmode}%
    {\scalebox{\scalefactor}{\DisplayProof}\proofSkipAmount \endcenter}
% ===========================================================

% TikZ
% ===========================================================
  \usepackage{tikz}
  \usepackage{tikz-cd}
  \usepackage{pgf}
  \usetikzlibrary{matrix}
  \usetikzlibrary{positioning}
  \usetikzlibrary{arrows}
  \usetikzlibrary{shapes}
  \tikzset{>=Latex}
  \usepackage[mode=image]{standalone}
% ===========================================================

% Seitenlayout, Kopf-/Fußzeile
% ===========================================================
  \usepackage{scrlayer-scrpage}
  \usepackage[top=3cm, bottom=2.5cm, left=2.5cm, right=2cm]{geometry}  
  \pagestyle{scrheadings}
  \clearscrheadfoot 
  \setheadsepline{0.4pt}
  \setfootsepline{0.4pt}
  \setkomafont{pagehead}{\normalfont\footnotesize}
  \setkomafont{pagefoot}{\normalfont\footnotesize}
  \lohead[]{\textit{Digitale Technologien -- Python Intro} -- WS 2026/2027}
  \rohead[]{%
    \ifthenelse{\boolean{loesungen}}{\textbf{Lösungen zu Blatt \blattnummer}}{Blatt \blattnummer}}
  %\rofoot[Version: \DTMnow]{%
  %  \ifthenelse{\boolean{loesungen}}{Version: \DTMnow}{}}
  %\lofoot[\jobname.tex]{%
  %  \ifthenelse{\boolean{loesungen}}{\jobname.tex}{}}
  \cfoot[\thepage]{\thepage}
  \raggedbottom
  \usepackage{setspace}          
  \linespread{1.2}
  \setlength{\parindent}{0pt}
  \setlength{\parskip}{0.5\baselineskip}
  \usepackage[all]{nowidow}
  \usepackage{lscape}
% ===========================================================

% Hyperref
% ===========================================================
  \usepackage[%
    hidelinks,
    pdfpagelabels,
    bookmarksopen=true,
    bookmarksnumbered=true,
    linkcolor=black,
    urlcolor=fbblau,
    plainpages=false,
    pagebackref,
    citecolor=black,
    hypertexnames=true,
    pdfborderstyle={/S/U},
    linkbordercolor=fbblau,
    colorlinks=false,
    backref=false]{hyperref}
  \hypersetup{final}
  \makeatletter
  \g@addto@macro{\UrlBreaks}{\UrlOrds}
  \makeatother
% ===========================================================

% Listen und Tabellen
% ===========================================================
  \usepackage{multicol}
  \usepackage{multirow}
  \usepackage[shortlabels]{enumitem}
  \setlist[enumerate]{font=\bfseries,itemsep=0.25\baselineskip}
  \setlist[itemize]{label=$\triangleright$,itemsep=-0.25\baselineskip}
  \usepackage{tabularx}
  \usepackage{array}
  \usepackage{listings}
  \lstset{breaklines=true,
    basicstyle=\ttfamily,
    frame=tblr,
    language=c,
    numbers=left}
  \usepackage{include/mips}

\definecolor{lightgreen}{rgb}{0,0.7,0}
\definecolor{gray}{rgb}{0.5,0.5,0.5}
\definecolor{mauve}{rgb}{0.58,0,0.82}

\lstdefinestyle{mips}{ %
  language=[mips]Assembler,       % the language of the code
  keywordstyle=\color{blue},          % keyword style
  keywordstyle=[2]\color{magenta},     % keyword style
  keywordstyle=[3]\color{red},          % keyword style
  commentstyle=\color{lightgreen},       % comment style
  stringstyle=\color{mauve},         % string literal style
  escapeinside={\%*}{*)},            % if you want to add a comment within your code
  morekeywords={*,...}               % if you want to add more keywords to the set
}

\newcommand{\mipsinline}[1]{\lstinline[style=mips,columns=fixed]{#1}}
\newcommand{\mipsregister}[1]{\mipsinline{\$#1}}
  \newcommand{\qed}{\hfill \ensuremath{\Box}}
  \newcommand{\cmd}[1]{\texttt{#1}}
% ===========================================================

% Sonstiges
% ===========================================================
  \usepackage[textsize=small,color=Red1!80!OrangeRed1!80]{todonotes}
  \usepackage{lipsum}
% ===========================================================

% Shortcuts and Math Commands
% ===========================================================
  \newcommand{\BB}{\mathbb{B}}      % Bool
  \newcommand{\CC}{\mathbb{C}}      % komplexe Zahlen
  \newcommand{\NN}{\mathbb{N}}      % nat. Zahlen
  \newcommand{\QQ}{\mathbb{Q}}      % rat. Zahlen
  \newcommand{\RR}{\mathbb{R}}      % reelle Zahlen
  \newcommand{\ZZ}{\mathbb{Z}}      % ganze Zahlen
  \newcommand{\oh}{\mathcal{O}}      % Landau-O
  \newcommand{\ol}[1]{\overline{#1}}    % überstreichen
  \newcommand{\wt}[1]{\widetilde{#1}}    % überschlängeln
  \newcommand{\wh}[1]{\widehat{#1}}    % überdachen
  \newcommand{\entspricht}{\mathrel{\widehat{=}}}
  
  \DeclarePairedDelimiter{\absolut}{\lvert}{\rvert}    % Betrag
  \DeclarePairedDelimiter{\ceiling}{\lceil}{\rceil}    % aufrunden
  \DeclarePairedDelimiter{\Floor}{\lfloor}{\rfloor}    % aufrunden
  \DeclarePairedDelimiter{\Norm}{\lVert}{\rVert}      % Norm
  \DeclarePairedDelimiter{\sprod}{\langle}{\rangle}    % spitze Klammern
  \DeclarePairedDelimiter{\enbrace}{(}{)}          % runde Klammern
  \DeclarePairedDelimiter{\benbrace}{\lbrack}{\rbrack}  % eckige Klammern
  \DeclarePairedDelimiter{\penbrace}{\{}{\}}        % geschweifte Klammern
  \newcommand{\Underbrace}[2]{{\underbrace{#1}_{#2}}}   % bessere Unterklammerungen
  
  % Kurzschreibweisen für Faule und Code-Vervollständigung
  \newcommand{\abs}[1]{\absolut*{#1}}
  \newcommand{\ceil}[1]{\ceiling*{#1}}
  \newcommand{\flo}[1]{\Floor*{#1}}
  \newcommand{\no}[1]{\Norm*{#1}}
  \newcommand{\sk}[1]{\sprod*{#1}}
  \newcommand{\enb}[1]{\enbrace*{#1}}
  \newcommand{\penb}[1]{\penbrace*{#1}}
  \newcommand{\benb}[1]{\benbrace*{#1}}
  \newcommand{\stack}[2]{\makebox[1cm][c]{$\stackrel{#1}{#2}$}}

% Übungszettel-Krams
% ===========================================================
  \usepackage[tikz]{mdframed}
  \newmdenv[%
    backgroundcolor = gray!20,
    linewidth=0pt,
    skipabove = 0.5cm,
    skipbelow = 0cm
  ]{notes}
  
  % Lösungsausgabe
  \newboolean{praesenzlsgn}
  \setboolean{praesenzlsgn}{\printpraesenzlsg}
  \newboolean{loesungen}
  \setboolean{loesungen}{\printloesungen}
  % Bewertungsschema-Ausgabe
  \newboolean{bewertungen}
  \setboolean{bewertungen}{\printbewertungen}

  \newcommand{\iforiginal}[1]{\ifthenelse{\boolean{praesenzlsgn}}{}{#1}}
  \newcommand{\ifpraesenzlsg}[1]{\ifthenelse{\boolean{praesenzlsgn}}{#1}{}}
  \newcommand{\ifloesung}[1]{\ifthenelse{\boolean{loesungen}}{#1}{}}
  \newcommand{\ifbewertung}[1]{\ifthenelse{\boolean{bewertungen}}{#1}{}}

  \newcommand{\umbruchoriginal}{\iforiginal{\newpage}}        % Seitenumbruch nur im Original
  \newcommand{\umbruchpraesenzlsg}{\ifpraesenzlsg{\newpage}}  % Seitenumbruch nur in Lösung d. Präsenzaufgabe
  \newcommand{\umbruchlsg}{\ifloesung{\newpage}}              % Seitenumbruch nur in Lösung
  \newcommand{\umbruchbewertung}{\ifbewertung{\newpage}}      % Seitenumbruch nur in Lösung mit Bewertungsschema
  \newcommand{\umbruchnurpraesenz}{\ifthenelse{\boolean{loesungen}}{}{\umbruchpraesenzlsg}}
  \newcommand{\umbruchnurlsg}{\ifthenelse{\boolean{bewertungen}}{}{\umbruchlsg}}

  \usepackage{comment}
  \newenvironment{praesenzlsg}%
    {\ifthenelse{\boolean{praesenzlsgn}}{\textbf{--- Präsenzaufgabe Anfang ---} \mbox{} \newline}{}}%
	{\ifthenelse{\boolean{praesenzlsgn}}{\mbox{} \newline \textbf{--- Präsenzaufgabe Ende ---}}{}}
  \newenvironment{loesung}%
    {\ifthenelse{\boolean{loesungen}}{\textbf{--- Lösung Anfang ---} \mbox{} \newline}{}}%
    {\ifthenelse{\boolean{loesungen}}{\mbox{} \newline \textbf{--- Lösung Ende ---} \ifthenelse{\boolean{bewertungen}}{}{\umbruchlsg}}{}}
  \newenvironment{bewertung}%
    {\ifthenelse{\boolean{bewertungen}}{\textbf{--- Bewertungsschema Anfang ---} \mbox{} \newline}{}}%
    {\ifthenelse{\boolean{bewertungen}}{\mbox{} \newline \textbf{--- Bewertungsschema Ende ---} \umbruchlsg}{}}

  \ifthenelse{\boolean{praesenzlsgn}}{}{\excludecomment{praesenzlsg}}
  \ifthenelse{\boolean{loesungen}}{}{\excludecomment{loesung}}
  \ifthenelse{\boolean{bewertungen}}{}{\excludecomment{bewertung}}

  % Aufgabennummerierung
  \newcounter{aufgabennummer} 
  \setcounter{aufgabennummer}{1}
  \newcommand{\aufgabe}[1]{\textbf{Aufgabe \blattnummer.\theaufgabennummer} \stepcounter{aufgabennummer} \hfill (#1 Punkte)}
  %\newcommand{\aufgabetitel}[2]{\iforiginal{\vspace{0.5cm}} \textbf{Aufgabe \blattnummer.\theaufgabennummer \ (#2)} \stepcounter{aufgabennummer} \hfill (#1 Punkte)}
  \newcommand{\aufgabetitel}[2]{\iforiginal{\vspace{0.5cm}} \textbf{Aufgabe \theaufgabennummer \ (#2)} \stepcounter{aufgabennummer} \hfill (#1 Punkte)}
  \newcommand{\praesenztitel}[1]{\iforiginal{\vspace{0.5cm}} \textbf{Aufgabe \blattnummer.\theaufgabennummer \ (#1)} \stepcounter{aufgabennummer} \hfill \smallpencil}

\newcommand{\danger}{\raisebox{0.5mm}{\fontencoding{U}\fontfamily{futs}\selectfont\char 66\relax}}
\newcommand{\ddanger}[1]{{\danger} \textbf{#1} {\danger}}

% Rechnerstrukturen-Krams
% ===========================================================
  \newcommand{\Band}{\ensuremath{\mathbin{\texttt{AND}}}}
  \newcommand{\Bor}{\ensuremath{\mathbin{\texttt{OR}}}}
  \newcommand{\Bnot}{\ensuremath{\mathbin{\texttt{NOT}}}}
  \newcommand{\Bnand}{\ensuremath{\mathbin{\texttt{NAND}}}}
  \newcommand{\Bnor}{\ensuremath{\mathbin{\texttt{NOR}}}}
  \newcommand{\Bxor}{\ensuremath{\mathbin{\texttt{XOR}}}}
  \newcommand{\Bxnor}{\ensuremath{\mathbin{\texttt{XNOR}}}}
  \newcommand{\Bone}{\ensuremath{\mathbin{\texttt{ONE}}}}
  \newcommand{\Bzero}{\ensuremath{\mathbin{\texttt{ZERO}}}}
  \newcommand{\Binary}[1]{\ensuremath{\texttt{#1}}\xspace}
  \newcommand{\ieee}[3]{\fbox{\ensuremath{\texttt{#1}}}\fbox{\ensuremath{\texttt{#2}}}\fbox{\ensuremath{\texttt{#3}}}}
  \DeclareMathOperator{\DNF}{DNF}
  \DeclareMathOperator{\KNF}{KNF}
  \DeclareMathOperator{\xor}{\oplus}
  \newcommand{\dc}{\star} %don't-care-Wert
 
\hypersetup{pdftitle={Digitale Technologien -- Übungsblatt 1: Python Intro},pdfauthor={Prof. Dr.-Ing. Sebastian Schlesinger}}
\lstset{language=Python,basicstyle=\small\ttfamily,numbers=none,showstringspaces=false,columns=fullflexible,keepspaces=true}
\begin{document}
\thispagestyle{scrplain}
\vspace*{-3cm}
\begin{minipage}[t][1.1cm][c]{4.5cm}
  \includegraphics[width=4cm]{include/hwr-logo.png}
\end{minipage}
\hfill
\begin{minipage}[t][1.5cm][c]{8cm}
  \begin{center}
  \begin{footnotesize}
    \textsf{Prof. Dr.-Ing. Sebastian Schlesinger} \\[-0.1cm]
    \textsf{Fachbereich 2 - Duales Studium Wirtschaft \& Technik}
    %\url{https://www.hwr-berlin.de}
  \end{footnotesize}
  \end{center}
\end{minipage}
\hfill
%\begin{minipage}[t][1.6cm][c]{3.2cm}
%  \includegraphics[width=2.8cm]{include/ESlogo2}
%\end{minipage}

\vspace*{-0.3cm}
\begin{center} 
  \hrulefill \\[0.1cm]
  {\large Übungsblatt } \blattnummer \\[0.15cm]
  {\huge \bfseries Digitale Technologien} \\[0.10cm]
  {Python Intro · Master Digitale Transformation} \\[0.10cm]
  {HWR Berlin, Wintersemester 2026/2027} \\[-0.4cm]
  \begin{tabular}{lcr}
    \hspace{0.3\textwidth}   & \hspace{0.3\textwidth} & \hspace{0.3\textwidth} \\
    Prof. Dr.-Ing. Sebastian Schlesinger   %&                        & "Ubung zur Klausurvorbereitung    \\ 
    %Blatt \blattnummer     % \abgabetermin          \\ 
  \end{tabular} \\[0.1cm]
  \hrulefill
\end{center}

\textbf{Lernziele:} Eine Python-Arbeitsumgebung einrichten, kleine Programme entwickeln,
Algorithmen nachvollziehen und Klassen sowie Polymorphie einsetzen.

\textbf{Arbeitsweise:} Aufgaben 1--6 bilden den Einstieg; 7--10 dienen der anschließenden
Vertiefung und dem Selbstlernen. Planen Sie insgesamt etwa 6--8 Stunden ein, abhängig
von Vorkenntnissen und Installationen. Das Begleitnotebook \texttt{Python\_Intro.ipynb}
(Abschnitte A--D und G) hilft bei der Bearbeitung. Insgesamt: \textbf{100 Übungspunkte}.
Schreiben Sie eigene Lösungen und erläutern Sie verwendete Hilfen kurz.

\aufgabetitel{12}{Python und VS Code einrichten}\\
\begin{enumerate}[a)]
\item Installieren Sie Python 3.12 oder neuer, VS Code sowie die Microsoft-Erweiterungen
\emph{Python} und \emph{Jupyter}. Prüfen Sie im Terminal die Python-Version:
\texttt{python3 --version} (macOS/Linux) bzw. \texttt{py -3 --version} (Windows).
\item Öffnen Sie einen neuen Projektordner \texttt{python-intro} in VS Code.
Erstellen Sie über \emph{Python: Create Environment} eine virtuelle Umgebung
\texttt{.venv} und wählen Sie sie über \emph{Python: Select Interpreter} aus.
\item Erstellen Sie \texttt{hello.py}, das Ihren Namen und \texttt{Hallo Python!} ausgibt.
Führen Sie die Datei aus. Geben Sie zusätzlich mit \texttt{sys.version} und
\texttt{sys.executable} Version und Pfad des verwendeten Interpreters aus.
\item Setzen Sie einen Breakpoint vor einer Ausgabe und starten Sie den Debugger.
Ändern Sie eine Variable und beobachten Sie ihren Wert. Erklären Sie in der
\texttt{README.md} den Unterschied zwischen Editor, Interpreter und virtueller Umgebung.
\end{enumerate}
\textbf{Nachweis:} \texttt{hello.py}, dokumentierte Python-Version, Interpreterpfad und kurze Erklärung.

\aufgabetitel{8}{Git und GitHub einrichten}\\
\begin{enumerate}[a)]
\item Erstellen Sie ein GitHub-Konto, installieren Sie Git und prüfen Sie \texttt{git --version}.
Legen Sie auf GitHub ein Repository \texttt{python-intro} mit README an und klonen Sie es
über VS Code (\emph{Git: Clone}) in einen separaten Ordner. Öffnen Sie diesen als Projekt.
\item Übernehmen Sie \texttt{hello.py} und Ihre README-Inhalte. Erstellen Sie die virtuelle
Umgebung im geklonten Projekt neu. Tragen Sie \texttt{.venv/}, \texttt{\_\_pycache\_\_/}
und \texttt{.ipynb\_checkpoints/} in \texttt{.gitignore} ein.
\item Konfigurieren Sie Ihren Git-Namen und Ihre Commit-E-Mail. Erstellen Sie einen Commit
und pushen Sie ihn. Legen Sie einen Branch \texttt{begruessung} an, ändern Sie die Begrüßung und pushen Sie den Branch.
Öffnen und mergen Sie einen Pull Request; aktualisieren Sie anschließend lokal \texttt{main}.
\item Erklären Sie Git vs. GitHub sowie Commit vs. Push. Halten Sie Repository-URL und
Commit-Kennung in der README fest. Zugangsdaten gehören nicht ins Repository.
\end{enumerate}
\textbf{Hilfe:} \href{https://code.visualstudio.com/docs/python/python-tutorial}{VS Code: Python-Tutorial};
\href{https://docs.github.com/en/get-started/using-github/hello-world}{GitHub: Hello World}.

\newpage
\aufgabetitel{6}{Google Colab nutzen}\\
\begin{enumerate}[a)]
\item Melden Sie sich bei \href{https://colab.research.google.com/}{Google Colab} an und erstellen
Sie ein Notebook \texttt{blatt01.ipynb}. Fügen Sie eine Markdown-Zelle mit Titel und Lernziel
sowie eine Codezelle mit \texttt{print("Hallo Colab!")} hinzu. Führen Sie diese aus.
\item Legen Sie in einer Zelle \texttt{preis = 19.90} und \texttt{menge = 3} fest und berechnen
Sie in einer zweiten Zelle den Umsatz. Erklären Sie, warum die Ausführungsreihenfolge zählt.
\item Starten Sie die Laufzeit neu und führen Sie alle Zellen von oben nach unten aus.
Laden Sie das Notebook als \texttt{.ipynb} herunter und nehmen Sie es in Ihr Git-Repository auf.
\item Nennen Sie je einen Vorteil der lokalen Arbeit und von Colab. Was müssen Sie sichern,
bevor eine Colab-Laufzeit endet? Unterscheiden Sie Notebook und Dateien im Laufzeitverzeichnis.
\end{enumerate}
\textbf{Nachweis:} Ausführbares Notebook und vier bis sechs Sätze in der README.
Die weiteren Programmieraufgaben können Sie lokal oder in Colab lösen.

\aufgabetitel{8}{Ein Python-Programm mit Docker ausführen}\\
\begin{enumerate}[a)]
\item Installieren und starten Sie Docker Desktop gemäß der
\href{https://docs.docker.com/get-started/get-docker/}{offiziellen Anleitung}
(unter Linux alternativ Docker Engine). Prüfen Sie im lokalen Terminal:
\begin{lstlisting}[language={}]
docker --version
docker run --rm hello-world
\end{lstlisting}
\item Erstellen Sie neben \texttt{hello.py} eine Datei \texttt{Dockerfile} mit folgendem Inhalt.
Für dieses Programm sind keine zusätzlichen Python-Pakete erforderlich.
\begin{lstlisting}[language={}]
FROM python:3.12-slim
WORKDIR /app
COPY hello.py .
CMD ["python", "hello.py"]
\end{lstlisting}
\item Bauen und starten Sie das Image im Projektordner:
\begin{lstlisting}[language={}]
docker build -t python-intro .
docker run --rm python-intro
\end{lstlisting}
Vergleichen Sie die Ausgabe mit der lokalen Ausführung. Ändern Sie die Begrüßung in
\texttt{hello.py}. Was wird ohne erneuten Build ausgegeben, was nach einem neuen Build?
\item Erklären Sie die Begriffe Dockerfile, Image und Container sowie den Unterschied
zwischen einem Container und einer Python-Umgebung \texttt{.venv}.
\end{enumerate}
\textbf{Nachweis:} Dockerfile, Befehle und beobachtete Ausgaben in der README.
Falls die Installation auf Ihrem Gerät scheitert, dokumentieren Sie Fehlermeldung und
Lösungsversuch; bearbeiten Sie die Python-Aufgaben zunächst in Colab weiter.

\begin{notes}
\textbf{Ergebnisse organisieren:} Speichern Sie Aufgaben 5--10 in \texttt{blatt01.py}
oder im Notebook. Funktionen sollen Ergebnisse mit \texttt{return} zurückgeben;
Ausgaben dienen nur der Darstellung. Ergänzen Sie die genannten Tests und mindestens
einen selbst gewählten Randfall pro Funktion. Nutzen Sie nur die Python-Standardbibliothek.
\end{notes}

\newpage
\aufgabetitel{12}{Kleine Programme: Variablen, Bedingungen und Schleifen}\\
\begin{enumerate}[a)]
\item \textbf{Temperaturumrechner (3 P.):} Lesen Sie eine Celsius-Temperatur mit \texttt{input}
ein, wandeln Sie sie mit \texttt{float} in eine Zahl um und berechnen Sie
$F=C\cdot 9/5+32$. Geben Sie das Ergebnis mit einer Nachkommastelle aus.
Prüfen Sie $0\mapsto32$, $100\mapsto212$ und $-40\mapsto-40$.
Warum ist eine Typumwandlung nötig?
\item \textbf{Versandkosten (3 P.):} Schreiben Sie \texttt{versand(warenwert)}.
Für einen nichtnegativen Warenwert unter 50 Euro betragen die Versandkosten 4{,}90 Euro,
ab 50 Euro sind sie null. Negative Werte sollen einen \texttt{ValueError} auslösen.
Prüfen Sie 0, 49{,}99, 50 und $-1$.
\item \textbf{FizzBuzz (3 P.):} Geben Sie die Zahlen 1 bis einschließlich 30 aus.
Ersetzen Sie durch 3 teilbare Zahlen durch \texttt{Fizz}, durch 5 teilbare durch
\texttt{Buzz} und durch beide teilbare durch \texttt{FizzBuzz}.
Nutzen Sie \texttt{\%}. Warum ist die Reihenfolge der Bedingungen wichtig?
\item \textbf{Umsätze (3 P.):} Gegeben sind \texttt{preise = [12, 8, 15, 10]} und
\texttt{mengen = [2, 5, 1, 3]}. Berechnen Sie mit einer Schleife die Einzelumsätze,
den Gesamtumsatz und die Anzahl der Einzelumsätze ab 30 Euro.
Erwartung: \texttt{[24, 40, 15, 30]}, Gesamtumsatz 109, Anzahl 2.
\end{enumerate}

\aufgabetitel{10}{Funktionen, Strings und Dictionaries}\\
\begin{enumerate}[a)]
\item \textbf{Durchschnitt (4 P.):} Implementieren Sie \texttt{durchschnitt(werte)} für eine
Liste von Zahlen ohne \texttt{sum}. Bei einer leeren Liste soll \texttt{ValueError}
ausgelöst werden. Prüfen Sie:
\begin{lstlisting}
assert durchschnitt([2, 4, 6]) == 4
assert durchschnitt([5]) == 5
assert durchschnitt([-4, 2]) == -1
\end{lstlisting}
Ergänzen Sie einen Test mit \texttt{try/except}, der den Fehlerfall überprüft.
\item \textbf{Palindrom (3 P.):} Implementieren Sie \texttt{ist\_palindrom(text)}.
Ignorieren Sie Groß-/Kleinschreibung und Leerzeichen; andere Zeichen bleiben erhalten.
\texttt{"Lagerregal"} und \texttt{"Ab ba"} liefern \texttt{True}, \texttt{"Python"}
liefert \texttt{False}. Der leere String gilt hier als Palindrom.
Tipp: \texttt{lower}, \texttt{replace} und Slicing.
\item \textbf{Häufigkeiten (3 P.):} Zählen Sie in
\texttt{["Web", "Shop", "Web", "Partner", "Web"]} die Häufigkeit jedes Werts
mit einem Dictionary und einer Schleife, ohne \texttt{Counter}.
Erwartung: \texttt{\{"Web": 3, "Shop": 1, "Partner": 1\}}.
Verpacken Sie den Code in \texttt{haeufigkeiten(werte)}; für \texttt{[]} ist das Ergebnis
\texttt{\{\}}. Tipp: \texttt{dict.get}.
\end{enumerate}
\textbf{Reflexion:} Erklären Sie an einer Ihrer Funktionen den Unterschied zwischen
\texttt{print} und \texttt{return}, zwischen Parameter und Argument sowie zwischen
\texttt{=} und \texttt{==}.

% \newpage
% \aufgabetitel{12}{Algorithmen: linear und binär suchen}\\
% Gegeben ist die aufsteigend sortierte Liste \texttt{werte = [3, 7, 11, 18, 23, 29, 42]}.
% Beide Suchfunktionen sollen den Index eines Treffers oder \texttt{-1} zurückgeben.
% \begin{enumerate}[a)]
% \item \textbf{Lineare Suche (4 P.):} Implementieren Sie \texttt{linear\_suchen(werte, ziel)}
% ohne \texttt{in} als Mitgliedschaftstest und ohne \texttt{index}.
% Durchlaufen Sie die Liste von links nach rechts. Bei mehrfachen Treffern soll der erste
% Index zurückgegeben werden. Diese Suche muss auch für unsortierte Listen funktionieren.
% \item \textbf{Binäre Suche (5 P.):} Implementieren Sie
% \texttt{binaer\_suchen(werte, ziel)} für aufsteigend sortierte Listen.
% Verwalten Sie eine linke und eine rechte inklusive Indexgrenze und prüfen Sie jeweils
% die Mitte \texttt{(links + rechts) // 2}. Verkleinern Sie den Suchbereich, bis ein Treffer
% gefunden oder der Bereich leer ist. Bei Duplikaten genügt ein beliebiger Trefferindex.
% \item \textbf{Vergleich (3 P.):} Notieren Sie für die Suche nach 29 die nacheinander
% geprüften Listenelemente beider Verfahren. Wie viele Listenelemente werden jeweils
% besucht? Erklären Sie die Sortierungsvoraussetzung und den Aufwand im schlechtesten
% Fall: $O(n)$ gegenüber $O(\log n)$. Warum lohnt sich vorheriges Sortieren nicht
% zwangsläufig für nur eine einzige Suche?
% \end{enumerate}
% \textbf{Tests für beide Verfahren:}
% \begin{lstlisting}
% werte = [3, 7, 11, 18, 23, 29, 42]
% # Jeweils beide Funktionsnamen einsetzen:
% assert linear_suchen(werte, 29) == 5
% assert linear_suchen(werte, 3) == 0
% assert linear_suchen(werte, 42) == 6
% assert linear_suchen(werte, 8) == -1
% assert linear_suchen([], 8) == -1
% \end{lstlisting}
% Prüfen Sie zusätzlich eine Liste mit nur einem Element sowie eine Liste mit Duplikaten.
% Die Eingabeliste soll unverändert bleiben.

% \aufgabetitel{10}{Algorithmen: Insertion Sort}\\
% Beim Sortieren durch Einfügen ist der linke Teil einer Liste bereits sortiert.
% Das nächste Element wird an der passenden Stelle in diesen Teil eingefügt.
% \begin{enumerate}[a)]
% \item \textbf{Von Hand (3 P.):} Beginnen Sie mit \texttt{[5, 2, 4, 1, 3]}.
% Notieren Sie den Zustand nach jedem der vier Einfügeschritte und markieren Sie den
% bereits sortierten Teil.
% \item \textbf{Implementierung (5 P.):} Schreiben Sie \texttt{insertion\_sort(werte)}.
% Arbeiten Sie auf einer Kopie der Eingabeliste und geben Sie diese sortiert zurück.
% Verschieben Sie größere Elemente mit einer \texttt{while}-Schleife nach rechts.
% Nutzen Sie weder \texttt{sort} noch \texttt{sorted} in Ihrer Implementierung.
% \item \textbf{Tests und Aufwand (2 P.):} Prüfen Sie \texttt{[]}, \texttt{[7]},
% \texttt{[3, 1, 3, -2]} sowie bereits sortierte und umgekehrt sortierte Listen.
% Prüfen Sie auch, dass die ursprüngliche Liste unverändert bleibt.
% In welchem Fall sind besonders viele Verschiebungen nötig? Begründen Sie den
% quadratischen Aufwand $O(n^2)$ im schlechtesten Fall.
% \end{enumerate}
% \textbf{Tipp:} Zum Überprüfen Ihrer Ergebnisse dürfen Sie \texttt{sorted} verwenden.

% \newpage
% \aufgabetitel{12}{Objektorientierung und Polymorphie}\\
% Modellieren Sie Produkte eines kleinen Onlineshops. Preise sind hier vereinfachend
% \texttt{float}-Werte in Euro; verwenden Sie bei Tests \texttt{math.isclose}.
% \begin{enumerate}[a)]
% \item \textbf{Klasse (4 P.):} Schreiben Sie \texttt{Produkt} mit
% \texttt{\_\_init\_\_(self, name, preis)} und den Instanzattributen \texttt{name} und
% \texttt{preis}. Negative Preise sollen bei der Erzeugung \texttt{ValueError} auslösen.
% Die Methode \texttt{endpreis(self)} liefert den Preis zurück.
% Erzeugen Sie zwei Objekte und erläutern Sie Klasse, Instanz, Attribut, Methode und \texttt{self}.
% \item \textbf{Vererbung (4 P.):} Leiten Sie \texttt{RabattProdukt} von \texttt{Produkt} ab.
% Der Konstruktor erhält zusätzlich \texttt{rabatt} als Zahl zwischen 0 und 1, etwa
% \texttt{0.2} für 20\,\%. Verwenden Sie \texttt{super()} und weisen Sie ungültige Rabatte
% mit \texttt{ValueError} zurück. Überschreiben Sie \texttt{endpreis} so, dass der Rabatt
% abgezogen wird. Testen Sie auch Rabatt 0 und 1.
% \item \textbf{Polymorphie (4 P.):} Schreiben Sie \texttt{gesamtpreis(produkte)}, das für
% jedes Objekt \texttt{endpreis()} aufruft und die Ergebnisse summiert.
% Verwenden Sie keine Typabfragen mit \texttt{type} oder \texttt{isinstance}.
% \begin{lstlisting}
% from math import isclose
% produkte = [Produkt("Kabel", 10),
%             RabattProdukt("Maus", 50, 0.2)]
% assert isclose(gesamtpreis(produkte), 50)
% assert isclose(gesamtpreis([]), 0)
% \end{lstlisting}
% Ergänzen Sie eine unabhängige Klasse \texttt{Gutschein} mit \texttt{endpreis()}, das
% \texttt{-5} liefert. Die gemischte Liste soll nun 45 ergeben, ohne
% \texttt{gesamtpreis} zu ändern. Erklären Sie daran Polymorphie und Duck Typing.
% \end{enumerate}

% \aufgabetitel{10}{Overriding und Overloading in Python}\\
% \begin{enumerate}[a)]
% \item \textbf{Gleicher Name (3 P.):} Sagen Sie voraus, was bei den beiden Aufrufen
% passiert. Führen Sie sie einzeln aus und erklären Sie das Ergebnis.
% \begin{lstlisting}
% def begruessen(name):
%     return "Hallo " + name

% def begruessen(vorname, nachname):
%     return "Hallo " + vorname + " " + nachname

% begruessen("Ada")
% begruessen("Ada", "Lovelace")
% \end{lstlisting}
% \item \textbf{Optionale Parameter (3 P.):} Ersetzen Sie die Definitionen durch eine
% Funktion \texttt{begruessen(vorname, nachname=None)}, die beide Aufrufe unterstützt.
% Erklären Sie, warum zwei gleichnamige Definitionen in Python keine automatische
% Auswahl anhand der Parameteranzahl bewirken. Grenzen Sie dies vom Überschreiben
% (Overriding) der Methode \texttt{endpreis} in Aufgabe 9 ab.
% \item \textbf{Operatorüberladung (4 P.):} Schreiben Sie \texttt{Vektor2D(x, y)} mit
% \texttt{\_\_add\_\_}, sodass \texttt{Vektor2D(1, 2) + Vektor2D(3, 4)} einen neuen Vektor
% mit \texttt{x == 4} und \texttt{y == 6} liefert. Die Operanden bleiben unverändert.
% Bei einem ungeeigneten Operandentyp geben Sie \texttt{NotImplemented} zurück.
% Testen Sie Addition, unveränderte Operanden und den \texttt{TypeError} bei Vektor plus Zahl.
% \end{enumerate}
% \textbf{Zum Abschluss:} Alle Dateien speichern, Skript bzw. Notebook von Beginn an ausführen,
% README mit Ausführungsanleitung ergänzen und den abschließenden Stand committen und pushen.
\end{document}
